%% ma: L11-B25-0006
%% lop: 11
%% chuong: 7
%% bai: 25
%% dang: 11.25.5
%% loai: TLN
%% muc: VD
%% dap_an: "0,65"
%% nguon: "(NBV)-ĐỀ SỐ 19-MỖI NGÀY 1 ĐỀ THI 2025.pdf (D19, MỖI NGÀY 1 ĐỀ - Đề số 19), Phần III Câu 2 (Chuyên Lương Thế Vinh – Đồng Nai 2025)"
%% kien_thuc: "Bài 25 (góc giữa hai mặt phẳng, góc nhị diện), hình chóp có mặt bên vuông góc với đáy"
%% trang_thai: chinh-thuc
%% ghi_chu: "Đáp án gốc 0,65 đúng (giải độc lập: cos = √21/7 ≈ 0,6547). Hình dựng lại theo hình gốc. Đã duyệt 19/10/2026."
\begin{ex}
Một công ty chế tạo tháp đang lên kế hoạch thiết kế một đài quan sát hình chóp $S.ABCD$. Đài quan sát có đáy là một sàn kính hình vuông với cạnh $a$, được đặt song song với mặt đất. Đỉnh $S$ của hình chóp là một cột ăng-ten phát sóng, được gắn cố định sao cho tam giác $SAB$ là tam giác đều và nằm trong một mặt phẳng vuông góc với đáy. Một camera giám sát được đặt trên cạnh $SC$ để ghi lại toàn cảnh đài quan sát, đồng thời một bảng điều khiển được đặt trong mặt phẳng $(SAB)$ để điều khiển camera. Hãy tính côsin của góc giữa hai mặt phẳng $(SCD)$ và $(SAB)$, giúp kỹ sư tính toán vị trí và góc quay tối ưu của camera (kết quả làm tròn đến hai chữ số phần thập phân).
\begin{center}
\begin{tikzpicture}[scale=.9]
  \path (0,0) coordinate (B) (3.5,0) coordinate (C) (4.7,1.3) coordinate (D) (1.2,1.3) coordinate (A);
  \coordinate (S) at (.6,3.9);
  \draw (S)--(B)--(C)--(D)--cycle (S)--(C);
  \draw[khuat] (S)--(A)--(B) (A)--(D);
  \path (S) node[above]{$S$} (A) node[left]{$A$} (B) node[below left]{$B$} (C) node[below right]{$C$} (D) node[right]{$D$};
\end{tikzpicture}
\end{center}
\shortans{0,65}
\loigiai{Gọi $H$, $K$ lần lượt là trung điểm của $AB$, $CD$. Vì $(SAB)\perp(ABCD)$ và $SAB$ đều nên $SH\perp(ABCD)$ và $SH=\dfrac{a\sqrt3}2$. Ta có $HK\perp AB$, $HK\perp SH$ nên $HK\perp(SAB)$, $HK=a$.

Trong tam giác $SHK$ kẻ $HI\perp SK$; vì $CD\perp(SHK)$ nên $HI\perp(SCD)$. Hai mặt phẳng $(SAB)$, $(SCD)$ có các đường vuông góc là $HK$ và $HI$ nên góc giữa chúng bằng góc $\widehat{IHK}$.

Trong tam giác vuông $SHK$: $\dfrac1{HI^2}=\dfrac1{SH^2}+\dfrac1{HK^2}=\dfrac4{3a^2}+\dfrac1{a^2}=\dfrac7{3a^2}$ nên $HI=\dfrac{a\sqrt3}{\sqrt7}$. Do đó $\cos\widehat{IHK}=\dfrac{HI}{HK}=\sqrt{\dfrac37}=\dfrac{\sqrt{21}}7\approx0{,}65$.}
\end{ex}
